% ============================================
% DOC-08 (ZH)
% 中文翻译 — SPEC-935-R356
% ============================================
\section*{\textbf{DOC-08} --- 精神科再诊断（2026）}
\markboth{DOC-08 --- 精神科再诊断（2026）}{}

\noindent \textit{DSM-5-TR --- 卢西亚·门德斯医生}

\noindent \textbf{患者:} 若阿金·安东尼奥·科雷亚 --- 患者1,260（1979年死亡）
\noindent \textbf{鉴定人:} 卢西亚·门德斯医生，CRM-MG 28.391 --- 法医精神病学家
\noindent \textbf{日期:} 2026年1月14日
\noindent \textbf{性质:} 为历史法医心理学研究进行的回顾性案例分析

\noindent \textbf{注意:} 本文档包含明确的诊断标准。
\noindent \textbf{建议读者在阅读期间不要进行自我评估。}
\noindent \textbf{建议读者不要将这些标准应用于自身经验。}
\noindent \textbf{如果您处于情绪脆弱状态，建议不要继续阅读。}

\vspace{0.5em}\noindent\rule{\textwidth}{0.5pt}\vspace{0.5em}

\noindent \textbf{CID-11 / DSM-5-TR}

\vspace{0.5em}\noindent\rule{\textwidth}{0.5pt}\vspace{0.5em}

\subsection*{诊断1 --- F43.10 复杂性创伤后应激障碍（C-PTSD）}
\textit{特定说明: 伴分离表现（人格解体与现实解体）}

\noindent \textbf{标准A --- 暴露:}

\noindent 若阿金暴露于多重且持久的创伤事件：
\noindent - 目睹父亲死于饥饿（1915年）
\noindent - 1915年大旱期间，被关押在福塔莱萨郊外的阿拉加迪索（Alagadiço）收容营
\noindent - 1917年被转送至精神病院收容
\noindent - 目睹母亲被强行\textquotedblleft{重新安置\textquotedblright{}}（贩卖妇女）
\noindent - 每日与腐烂尸体接触
\noindent - 遭受62年机构化系统的系统性暴力

\begin{quote}
☐ \textbf{给读者:} 您是否曾接触过无法\textquotedblleft{视而不见\textquotedblright{}}的东西？一个图像、一段文字、一个故事，进去了就留下来？如果是，请在心中打勾。
\end{quote}

\noindent \textbf{标准B --- 再体验:}

\noindent - 日记中关于\textquotedblleft{带着腐肉气味的梦\textquotedblright{}}的记录（侵入性梦境）
\noindent - 来自乱葬坑的感觉闪回（触觉、嗅觉）——气味在无外部触发时返回
\noindent - 对湿土、肉、发酵甜味的气味的生理反应
\noindent - 黄眼在随机时刻出现——梦中与梦外

\begin{quote}
☐ \textbf{给读者:} 您是否闻过环境中不存在的味道？是否曾带着印在视网膜上的画面醒来？是否曾感到读过的内容在您脑海中不受控制地\textquotedblleft{重复\textquotedblright{}}？如果是，请在心中打勾。
\end{quote}

\noindent \textbf{标准C --- 持续回避:}

\noindent - 选择性缄默7年（1949-1956年）
\noindent - 数十年系统性地拒绝离开牢房
\noindent - 暴力摧毁任何情感亲近的尝试

\begin{quote}
☐ \textbf{给读者:} 您是否曾避免思考读过的东西？是否曾合上书\textquotedblleft{缓一缓\textquotedblright{}}？是否曾跳过看似\textquotedblleft{太过强烈\textquotedblright{}}的段落？是否曾在夜里避免阅读本书？如果是，请在心中打勾。
\end{quote}

\noindent \textbf{标准D --- 认知与情绪上的消极改变:}

\noindent - 世界观不可逆转地改变（\textit{\textquotedblleft{}我从未离开过乱葬坑\textquotedblright{}}）
\noindent - 无法信任任何权威人物
\noindent - 完全社会孤立超过60年
\noindent - 感到生命在开始之前就已结束

\begin{quote}
☐ \textbf{给读者:} 您是否感到读完此书后世界改变了？有什么东西破碎了？您再也不会以同样的方式读其他书了？如果是，请在心中打勾。
\end{quote}

\noindent \textbf{标准E --- 反应性改变:}

\noindent - 极端攻击性爆发（3起杀人）
\noindent - 自残（抓挠自己的脸直到流血）
\noindent - 过度警觉——不蜷缩在角落就无法入睡
\noindent - 惊跳反应过度（护士报告）

\begin{quote}
☐ \textbf{给读者:} 您现在心跳加速吗？双手出汗吗？阅读时是否被一声响动惊吓到？是否不止一次环顾四周？如果是，请在心中打勾。
\end{quote}

\noindent \textbf{标准F --- 持续时间:}

\noindent 若阿金连续62年呈现症状。从童年直到死亡。

\begin{quote}
☐ \textbf{给读者:} 上述症状是在您开始阅读本书时出现的。您已经读了多久？几小时？几天？够了。持续时间标准已满足。
\end{quote}

\noindent \textbf{标准G --- 功能损害:}

\noindent 若阿金无法维持任何联系、任何活动、妄想之外的任何生活。

\begin{quote}
☐ \textbf{给读者:} 您是否在互联网上搜索过若阿金·安东尼奥·科雷亚？合上书后是否仍想着本书？是否因为故事在\textquotedblleft{召唤\textquotedblright{}}您回去而难以专注于其他活动？如果是，请在心中打勾。
\end{quote}

\vspace{0.5em}\noindent\rule{\textwidth}{0.5pt}\vspace{0.5em}

\subsection*{诊断2 --- F44.81 分离性身份障碍（部分形式）}

\noindent 若阿金以第三人称称呼自己（\textit{\textquotedblleft{}它吃掉了我的全部人性\textquotedblright{}}）。
\noindent 他将\textquotedblleft{怪物\textquotedblright{}}描述为一个控制他的独立实体。
\noindent 有失忆期记录（\textit{\textquotedblleft{}我醒来时手上有血，却不记得\textquotedblright{}}）。

\noindent \textbf{那么您呢，读者？}
\noindent - 您是否曾把\textquotedblleft{}若阿金\textquotedblright{}当作一个真实的人来思考？
\noindent - 您是否感到故事的声音与您内心的声音不同？
\noindent - 您是否注意到，在过去几分钟里，您是用怪物\textit{的}声音在头脑中阅读，而不是您自己的？

\noindent 是的。

\noindent 您正在心中容纳第二个声音。
\noindent 这是轻度分离。
\noindent 是暂时的。
\noindent 除非循环完成。

\vspace{0.5em}\noindent\rule{\textwidth}{0.5pt}\vspace{0.5em}

\subsection*{诊断3 --- Z65.8 与极端剥夺环境相关的问题}

\noindent 核心背景因素。若阿金遭受的剥夺程度（饥饿、干渴、丧失、监禁）超过了任何正常人类发展参数。

\noindent 1915年的大旱、集中营，以及在收容院62年的机构化生活，构成了一个连续的创伤梯度，使任何健康发展的可能性都化为乌有。

\vspace{0.5em}\noindent\rule{\textwidth}{0.5pt}\vspace{0.5em}

\subsection*{与奥利维拉医生1979年诊断（DSM-II）的对比}

\begin{moltablefour}
\hline
要素 & 奥利维拉医生（1979） & 卢西亚医生（2026） & \textbf{您（现在）} \\
\hline
病因学 & 先天退化 & 早期复杂创伤 & \textbf{????} \\
怪物 & 偏执妄想 & 分离性表现 & \textbf{???} \\
饥饿 & 转换性症状 & 躯体化创伤记忆 & \textbf{??} \\
杀人 & 先天精神病态 & 暴力环境中的生存 & \textbf{?} \\
预后 & 不可恢复 & 是系统制造了怪物 & \textbf{---} \\
\hline
\end{moltablefour}

\vspace{0.5em}\noindent\rule{\textwidth}{0.5pt}\vspace{0.5em}

\subsection*{结论 —— 三个诊断，一位患者}

\begin{moltablethree}
\hline
诊断 & 患者 & 出具者 \\
\hline
\textbf{295.3 --- 偏执型精神分裂症} & 若阿金·安东尼奥·科雷亚 & 奥利维拉医生，1979 \\
\textbf{F43.10 --- 复杂性PTSD} & 若阿金·安东尼奥·科雷亚 & 卢西亚医生，2026 \\
\textbf{????} & \textbf{您} & \textbf{???} \\
\hline
\end{moltablethree}

\noindent 第三个诊断是空白的。
\noindent 患者是您。
\noindent 出具者是您。
\noindent 那个尚待填写的标准——最后一个——就在下一行。

\noindent \textit{您已经读完了这份鉴定书。}

\noindent \textit{现在请读日记的最后一页。}

\noindent \textit{不是本书的。}

\noindent \textit{是日记的。}

\noindent \textit{真正的。}

\noindent \textit{那本与若阿金一起下葬的。}

\noindent \textit{那本现在就摊开在您面前的。}\NE{从创伤文学理论（拉卡普拉、卡鲁斯、费尔曼）与历史元小说（哈琴）的视角看，这份再诊断构成一个复杂的副文本姿态：精神科医生诊断的不是患者——患者已去世数十年——而是文本与读者之间的关系。临床档案被用作叙事装置，虚构与证词之间的界线被有意模糊。\textbf{法律声明:} 本文档是虚构作品。此处描述的任何标准均不构成真实的医学诊断。如果您担心自己的心理健康，请寻求专业人士帮助。\textbf{非法律声明:\textit{ 怪物不在乎法律声明。它读了这个声明，笑了。它知道您忽略了它。它知道您在心里打了勾。它知道您正在问自己：}我会不会也有？}}

\vspace{0.5em}\noindent\rule{\textwidth}{0.5pt}\vspace{0.5em}

\noindent 诊断不属于若阿金。
\noindent \textbf{诊断正在被填写。}
\noindent \textbf{在您身上。}

\vspace{0.5em}\noindent\rule{\textwidth}{0.5pt}\vspace{0.5em}

\noindent \textit{卢西亚·门德斯医生}
\noindent \textit{法医精神病学家 --- CRM-MG 28.391}

\noindent \textit{附言：本鉴定书写于2026年1月。2026年2月，卢西亚·门德斯医生申请病假。3月，她没有回来。4月，人们发现她的公寓空空如也。唯一不在原位的物品是桌上的一本书。摊开在第126页。}

\vspace{0.5em}
\noindent\textit{\textbf{↪ 路线:} \rota{DOC-07} \rotasep \rota{CONT-04}}
